% ============================================================================
% PARKED MATERIAL --- NOT FOR SUBMISSION
%
% Included from main.tex behind \iftoggle{scratch}. Set the toggle to false
% (or delete the \input line) to drop this appendix before submission.
% ============================================================================

\section{Scratch}
\label{app:scratch}

Material cut from the body during the 2026-09 restructure and parked here
pending a decision. \textbf{This appendix is not part of the paper.}

\subsection{Contributions list (cut from the introduction)}

The introduction previously closed with an enumerated contributions list. It
was cut because the section it headed is now prose; the claims themselves are
made and evidenced in the body.

\begin{enumerate}
	\setlength{\itemsep}{1pt}
	\item \textbf{The claim ladder} (\cref{tab:ladder}) as an organising frame
	for passive compute verification: four claims, their information costs, and
	their ceilings, closed out one by one in \cref{tab:closeout}.
	\item \textbf{A tracked cyclostationary detector for Rung~1}
	(\cref{sec:methods}), with a negative methodological result we report
	rather than bury: no covariance configuration calibrates the complete
	tracker--resample--test pipeline end to end, so it is not a constant
	false-alarm-rate test and we do not present it as one. It is an effective
	score with per-trace surrogate calibration, level-exact on the
	linear-stationary nulls we tested. A pre-registered bake-off
	(\cref{app:bakeoff}) makes the choice of detector an outcome, and against a
	Whittle likelihood-ratio ceiling the tracker is essentially optimal on the
	inference null while only the cyclostationary order family approaches the
	ceiling on structural confusers.
	\item \textbf{The de-periodicisation frontier} (\cref{sec:frontier}), the
	central result, priced against measured throughput anchors. Its shape is an
	asymmetry: against fixed-frequency tests a prover reaches a hiding rate of
	$0.74$ at essentially zero throughput cost by varying real work per
	iteration, whereas against the tracking class the best anchored attack
	reaches only $0.16$ at $159\%$ overhead.
	\item \textbf{A minimum meter specification} (\cref{subsec:meter_spec}):
	the observation channel is a first-class condition on every rung, and we
	locate where each detector class dies in sample rate, integration window,
	and in-band transfer function. A $1$\,s-integrating, $1$\,Hz-reporting meter
	defeats every detector on an unmodified honest workload.
	\item \textbf{Rung-2 classification with semantic falsification controls}
	(\cref{subsec:rung2}). Three decision rules separate training from the
	stated inference null perfectly and survive six domain shifts; but five
	non-training loads --- three of them benign operational loads, alongside a
	decoy that computes gradients and discards them --- all score as training. The
	meter certifies the physical schedule, not training semantics: a ceiling no
	better passive detector removes, and one that limits positive evidence as much
	as it admits counterfeits.
	\item \textbf{A negative transport case} (\cref{sec:reality}): what we
	measured on real hardware, why it does not carry the distributed claim, and
	what would falsify the prediction that it should.
\end{enumerate}

Supporting these, \cref{subsec:identifiability} gives the formal counterpart of
the Rung-1 and Rung-2 ceilings: an observational-equivalence argument for why
no passive test can identify training semantics, and a covertness cost bound
showing that a prover using i.i.d.\ phase jitter to reach a target covertness
level against an exhibited verifier must pay in throughput or learning
efficiency. We are explicit about what that bound does not cover, and
\cref{sec:frontier} supplies the near-free attack that shows why the gap
matters.
